% Generated draft: requires mathematical and source review.
\documentclass[11pt]{article}
\usepackage[margin=25mm]{geometry}
\usepackage{fontspec}
\setmainfont{Latin Modern Roman}
\usepackage{amsmath,amssymb,mathtools,unicode-math}
\setmathfont{Latin Modern Math}
\usepackage{enumitem,needspace,xurl,hyperref}
\setlength{\parindent}{0pt}
\setlength{\parskip}{5pt}
\setlength{\emergencystretch}{3em}
\newcommand{\N}{\mathbb{N}}
\newcommand{\Z}{\mathbb{Z}}
\newcommand{\Q}{\mathbb{Q}}
\newcommand{\R}{\mathbb{R}}
\newcommand{\C}{\mathbb{C}}
\newcommand{\F}{\mathbb{F}}
\newcommand{\A}{\mathbb{A}}
\newcommand{\PP}{\mathbb{P}}
\newcommand{\E}{\mathbb{E}}
\newcommand{\eps}{\varepsilon}
\newcommand{\dd}{\,\mathrm{d}}
\begin{document}
\section*{UnsolvedMath ID: 20003164}
\textbf{Source problem code:} catalogue entry
\subsection{Definitions and mathematical statement}
A polar action is an isometric action of a Lie group \(G\) on a Riemannian manifold \(M\) admitting a complete section \(\Sigma\) meeting every orbit orthogonally. Depending on convention a section may be immersed rather than embedded. The source asks whether rational ellipticity of \(\Sigma\) forces rational ellipticity of \(M\). The following extracted heading ``5. Vector Bundles'' belongs to the next section and is excluded.

In the usual simply connected finite-CW setting, rational ellipticity means finite-dimensional rational cohomology and finite total rational homotopy. For spaces outside that setting, nilpotence, fundamental-group and finite-type conventions must be specified. To state the supplied calculation unambiguously, put
\[
e_\Q(Z)=\sum_{i\ge2}\dim_\Q(\pi_i(Z)\otimes\Q),
\]
allowing \(+\infty\). This higher-homotopy rank alone is not silently declared the complete definition of rational ellipticity for arbitrary nonnilpotent spaces.

The aid treats a compact Lie group acting polarly on a closed connected manifold without singular orbits, meaning no orbit has smaller dimension than the principal orbit. Exceptional orbits of the same dimension are permitted. Let \(H\) be principal isotropy. The reported exceptional-polar reconstruction is a quotient of \(G/H\times\Sigma\) by the appropriate Weyl/holonomy action, giving the product covering used in the argument. The aid calls this a finite covering; properness/compactness of the section and finiteness of the Weyl action must be checked, since those do not follow merely by calling an arbitrary section immersed. For the stated compact exceptional reconstruction with the required covering hypotheses,
\[
\pi_i(M)\otimes\Q\cong
(\pi_i(G/H)\otimes\Q)\oplus(\pi_i(\Sigma)\otimes\Q),\qquad i\ge2,
\]
and hence
\[
e_\Q(M)=e_\Q(G/H)+e_\Q(\Sigma),\qquad e_\Q(G/H)<\infty.
\]
Coverings preserve higher homotopy groups; the splitting itself requires a product covering, not that its degree be finite. The compact homogeneous orbit has finite total higher rational homotopy rank by the compact-Lie-group fibration.

Thus finiteness or infinitude of higher rational homotopy rank is equivalent for \(M\) and \(\Sigma\) in this reconstruction class. When the standard rational-ellipticity hypotheses are also satisfied, this proves the requested implication there. The aid consequently localizes any compact counterexample within its intended reconstruction setting to an action with singular isotropy. It does not establish the implication for arbitrary singular polar actions, unspecified section conventions or every extension of rational ellipticity beyond nilpotent finite-type spaces. The original problem remains open in the supplied account.
\subsection{Short English statement}
Ask whether rational ellipticity of a polar section transfers to the ambient manifold. The supplied no-singular-orbit calculation splits higher rational homotopy through a product covering; section properness and the usual rational-homotopy hypotheses must be checked before interpreting it as the full ellipticity assertion.
\subsection{Sources}
\begin{itemize}
\item Problem 39, AIM workshop \emph{Manifolds with nonnegative sectional curvature}, \url{https://aimath.org/WWN/nnsectcurvature/nnsectcurvature.pdf}. The supplied verification removes the next-section heading from the question.
\item Local UnsolvedMath registry, permanent ID 20003164, exact code \texttt{catalogue entry}; canonical location \texttt{aim-topology-notes.json}, notes[251], dated 14 August 2026; statement labeled exact.
\item Research attempt 1: partial result, problem open at the run, rigor label proved, novelty confidence low; registry status partially solved. The exceptional-polar reconstruction and rank formula are supplied machine-generated research aids, not peer reviewed; independent verification is required before citation. No full external reconstruction or rational-homotopy bibliography is supplied.
\end{itemize}
\end{document}
